\section{Antrag auf Anschubförderung --- UnRAVeL}

\noindent Sehr geehrte Leiter des Graduiertenkollegs UnRAVeL,

hiermit beantrage ich für den Zeitraum vom 01.10.2026 bis zum 31.12.2026 eine Postdoc-Anschubförderung (Start-up Funding) im Graduiertenkolleg UnRAVeL.

Seit Oktober 2022 bin ich als Doktorand und wissenschaftlicher Mitarbeiter am Lehrstuhl für Programmiersprachen und Verifikation der RWTH Aachen tätig und seit dieser Zeit zugleich Stipendiat von UnRAVeL.

In meiner Forschung entwickle ich automatische Verfahren, mit denen sich das Verhalten von Computerprogrammen mathematisch überprüfen lässt.
Im Mittelpunkt stehen dabei zwei grundlegende Fragen der Informatik: Hält ein Programm garantiert an (\emph{Terminierung}), und wie lange läuft es (\emph{Laufzeit})?
Im Rahmen von UnRAVeL analysiere ich \emph{probabilistische Programme}, die während ihrer Ausführung Zufallsentscheidungen treffen können.
Wahrscheinlichkeiten sind für viele moderne Anwendungen zentral, beispielsweise für randomisierte Algorithmen, die Kryptographie oder das maschinelle Lernen.
Um derartige Programme zu analysieren, verwende ich \emph{Termersetzungssysteme}: ein einfaches, aber ausdrucksstarkes Berechnungsmodell, das sich besonders gut eignet, um Programmen mit Datenstrukturen zu analysieren.

Meine Dissertation mit dem Titel ``Automatically Analyzing Termination and Expected Runtime Complexity of Probabilistic Term Rewriting'' werde ich im Juli 2026 einreichen.
Sie wird von Prof.\ Dr.\ rer.\ nat.\ Jürgen Giesl (RWTH Aachen University) als Erstgutachter und von Prof.\ Dr.\ Ugo Dal Lago (Universität Bologna) als Zweitgutachter betreut.
Die Verteidigung findet voraussichtlich im September oder Oktober 2026 statt.

Während meiner Promotion habe ich gelernt, eigenständig wissenschaftlich hochwertige Arbeiten zu entwickeln und zu veröffentlichen.
Dies spiegelt sich in meiner Publikationsliste wider, in der ich überwiegend als Erstautor auftrete.
Mein Betreuer gab dabei die allgemeine Forschungsrichtung vor (die Analyse probabilistischer Termersetzung). Die konkreten Ansätze, um diese Analyse zu verbessern, habe ich stets selbstständig erarbeitet.
So konnte ich meine im probabilistischen Kontext entwickelten Methoden sogar nutzen, um ein seit Jahrzehnten offenes Problem der klassischen (nicht-probabilistischen) Termersetzung zu lösen: ein \emph{Dependency Pair Framework} für die relative Terminierung von Termersetzungssystemen.
Die Anschubförderung würde es mir ermöglichen, diese Selbstständigkeit weiter auszubauen und einen eigenen Forschungsantrag zur Einwerbung von Drittmitteln vorzubereiten.

Für den Förderzeitraum plane ich, die folgenden Projekte abzuschließen und fortzuführen:

\begin{itemize}
    \item \emph{Decidability Results for Almost-Sure Termination of Probabilistic Term Rewriting}: Aufbauend auf einer ausgezeichneten Bachelorarbeit von Arion Scheid möchte ich untersuchen, für welche Klassen von Programmen sich die Terminierung automatisch entscheiden lässt. Ein Konferenzpapier ist für Ende 2026 geplant.
    \item \emph{Reachability Analysis in Probabilistic Term Rewriting}: Meine Dissertation betrachtet ausschließlich die Terminierung. Mindestens ebenso wichtig für die Software-Verifikation ist jedoch die Frage nach der Sicherheit eines Programms: Kann das Programm einen unsicheren Zustand erreichen? Hierzu möchte ich erste Resultate zur Erreichbarkeitsanalyse für probabilistische Termersetzung erzielen. Auch hier ist ein Konferenzpapier für Ende 2026 geplant.
    % \item \emph{Analyzing Termination of Probabilistic Integer Term Rewriting}: Bislang beschränken sich meine Verfahren auf Programme über Datenstrukturen. Diese lassen sich damit zwar gut beschreiben, eignen sich aber schlecht für das Rechnen mit Zahlen. Indem ich mein Modell mit \emph{Integer-Transitionssystemen} kombiniere, die ganze Zahlen direkt unterstützen, möchte ich die Analyse realistischerer Programme ermöglichen.
    \item \emph{Confluence Analysis of Weighted Rewriting} mit Emma Ahrens (Gruppe von Prof. Katoen): Die von uns entwickelte ``gewichtete'' Semantik für Ersetzungssysteme möchte ich um eine Analyse der \emph{Konfluenz} erweitern: liefert ein Program unabhängig von der Reihenfolge seiner Rechenschritte stets dasselbe Ergebnis und verursacht es dabei stets die selben Kosten?
    \item \emph{How to Nicely Integrate Formal Verification}: Gemeinsam mit Prof.\ Tom Beckmann (Kyoto University of Advanced Science) möchte ich der bereits während eines Forschungsaufenthalts am HPI Potsdam begonnenen Frage nachgehen, wie sich formale Verifikation möglichst einfach in den Software-Entwicklungsprozess integrieren lässt.
\end{itemize}

Darüber hinaus plane ich, wie bereits erwähnt, im Rahmen der Anschubförderung einen Forschungsantrag bei der Deutschen Forschungsgemeinschaft (DFG) vorzubereiten.
Konkret strebe ich einen Antrag im \emph{Walter-Benjamin-Programm} der DFG an.
Die oben genannten Projekte bilden die inhaltliche Grundlage dieses Antrags.
% , und die Anschubförderung verschafft mir den nötigen zeitlichen Freiraum, ihn im Anschluss an die Einreichung meiner Dissertation sorgfältig auszuarbeiten.

% Mein zentraler wissenschaftlicher Beitrag ist das \textit{Annotated-Dependency-Pair-(ADP)-Framework} zum Beweisen und Widerlegen fast-sicherer Terminierung sowie zur Beschränkung erwarteter Laufzeiten. Diese Methoden sind im Werkzeug AProVE implementiert, das ich seit 2023 als Hauptentwickler betreue und das bei der jährlichen International Termination and Complexity Competition (termCOMP) von 2023 bis 2025 in zahlreichen Kategorien den ersten Platz belegte. Meine Ergebnisse habe ich bislang in 16 Publikationen veröffentlicht, darunter Beiträge auf führenden Konferenzen wie CADE, FoSSaCS, FLOPS, IJCAR und PPDP sowie in den Zeitschriften \textit{Science of Computer Programming} und \textit{Logical Methods in Computer Science}.\par

% Mein Informatikstudium an der RWTH Aachen habe ich 2022 mit Auszeichnung (Note~1{,}2) abgeschlossen und wurde dafür mit der Springorum-Denkmünze ausgezeichnet. Neben der Forschung bin ich aktiv in der Lehre tätig: Ich war unter anderem als Übungsleiter für die Vorlesungen ``Grundlagen der funktionalen Programmierung'', ``Grundlagen der logischen Programmierung'' und ``Programmierung'' verantwortlich --- jeweils verbunden mit einer Nominierung für den Lehrpreis --- und habe bisher 24 Abschlussarbeiten betreut. Dem Graduiertenkolleg UnRAVeL bin ich darüber hinaus eng verbunden: Ich habe sowohl das UnRAVeL-Symposium 2026 als auch den UnRAVeL-Herbstworkshop 2023 organisiert.\par

% \softtext{[Hier folgt die konkrete Beschreibung des Vorhabens, für das die Anschubförderung beantragt wird: Zielsetzung und wissenschaftliche Motivation, geplante Aktivitäten (z.\,B.\ Forschungsaufenthalt, Kooperation oder Workshop), die benötigten Mittel sowie der Bezug zu den Forschungszielen von UnRAVeL.]}\par

Über eine positive Rückmeldung zu meinem Antrag und die weitere Zusammenarbeit im Graduiertenkolleg UnRAVeL würde ich mich sehr freuen.
Für Rückfragen stehe ich Ihnen jederzeit gerne zur Verfügung.
